\documentclass[conference]{IEEEtran}
\usepackage[T1]{fontenc}
\usepackage{graphicx,booktabs,array,calc,amsmath,url}
\usepackage[hidelinks]{hyperref}
\newcommand{\real}[1]{#1}
\title{Design and Evaluation of a Resilient Private Network Service Platform}
\author{\IEEEauthorblockN{Saksham Narotra, Shrihari K.N, Rachit Singh, and Amogha Raj Sandur}
\IEEEauthorblockA{Team 1.1.1.1 --- Computer Networks Course Project --- November 2026}}
\begin{document}
\maketitle
\begin{abstract}
This report presents the design,
implementation and evaluation of a fully local private network service
platform using four macOS laptops. Phase 1 established LAN connectivity,
private DNS resolution, two Python HTTP backends, an nginx reverse proxy
with round-robin load balancing, trusted HTTPS, conditional HTTP cache
validation, and packet-level observation. Phase 2 introduced redundant
DNS resolvers, a 30 s DNS time-to-live (TTL), PF-based backend
isolation, passive upstream failover, controlled migration to a standby
TLS edge, and systematic fault diagnosis. Recorded results show
continued application availability during a primary DNS outage, ten
successful requests through the remaining backend during a controlled
backend failure, and successful HTTPS cutover and restoration. The
project demonstrates the distinction between name resolution, transport
connectivity, TLS trust and application health. Limitations include
manual DNS cutover, non-permanent firewall isolation, and incomplete
contemporaneous evidence for some experiments.
\end{abstract}
\begin{IEEEkeywords}
DNS, HTTP, TLS, TCP, reverse proxy, load
balancing, high availability, network isolation, Wireshark.
\end{IEEEkeywords}
\section{INTRODUCTION}\label{i.-introduction}

Networked applications depend on a sequence of independent services and
protocols. A hostname must resolve to an address; a transport connection
must reach the correct socket; a secure session must be established; and
the application must return a valid response. A failure at any layer can
make a working application appear unavailable. The course project
\cite{ref1}  provides a local setting in which these dependencies can be
configured, observed and deliberately disrupted without cloud
infrastructure.

The objective of Team 1.1.1.1 was to implement an end-to-end service
using a private domain, app.team1.test, and to increase resilience
without rebuilding the Phase 1 platform. A second name, api.team1.test,
resolves to the same edge. The application is deliberately minimal,
allowing attention to remain on addressing, DNS, transport, security,
HTTP and service recovery. All runtime components were hosted on the
team laptops; GitHub was used for version control and evidence storage
rather than hosting \cite{ref2}.

The contributions of this work are: (1) a reproducible four-node
topology with trusted name-based access; (2) an observation of DNS, TCP,
TLS and upstream HTTP through packet evidence; and (3) controlled
experiments addressing resolver, backend and edge failures, including
rollback procedures. The report evaluates the recorded experiments
rather than presenting them as production-grade availability benchmarks.

\section{SYSTEM ARCHITECTURE}\label{ii.-system-architecture}

\subsection{Physical and Logical
Topology}\label{a.-physical-and-logical-topology}

Four macOS laptops were connected to a private wireless LAN. Phase 1
used a lab WLAN; Phase 2 used a controlled hotspot with interface en0
and gateway 10.185.128.178. DHCP supplied IPv4 addresses dynamically, so
the same logical roles appear with different IPs in earlier screenshots.
The baseline Phase 1 capture used Mac 1 at 10.7.1.82, Mac 2 at
10.7.20.19, Mac 3 at 10.7.13.26, and Mac 4 at 10.7.1.23 \cite{ref2}. Table I
shows the last recorded Phase 2 addresses. Configuration files must be
adjusted if those addresses change.

\begin{figure}[!t]\centering
\includegraphics[width=\columnwidth]{ieee_tex_assets/media/image1.png}
\caption{Phase 2 topology. DNS returns an edge address; the client then
connects directly to the edge. The standby role is activated only during
migration.}
\end{figure}

\begin{table}[!t]
\caption{PHASE 2 NODE AND SERVICE INVENTORY}
\centering\footnotesize
\begin{tabular}{@{}
  >{\raggedright\arraybackslash}p{(\columnwidth - 4\tabcolsep) * \real{0.3070}}
  >{\raggedright\arraybackslash}p{(\columnwidth - 4\tabcolsep) * \real{0.3011}}
  >{\raggedright\arraybackslash}p{(\columnwidth - 4\tabcolsep) * \real{0.3919}}@{}}
\toprule
\begin{minipage}[b]{\linewidth}\raggedright
\textbf{Node / member}
\end{minipage} & \begin{minipage}[b]{\linewidth}\raggedright
\textbf{IPv4 address}
\end{minipage} & \begin{minipage}[b]{\linewidth}\raggedright
\textbf{Role and port}
\end{minipage} \\
\midrule



Mac 1 / Saksham & 10.185.128.177 & Primary DNS, 53; test client \\
Mac 2 / Shrihari & 10.185.128.147 & nginx, 8081/8443 \\
Mac 3 / Rachit & 10.185.128.163 & Backend A, 3001; standby 8443 \\
Mac 4 / Amogha & 10.185.128.32 & Backend B, 3002; backup DNS, 53 \\
\bottomrule
\end{tabular}
\end{table}

\subsection{End-to-End Request
Processing}\label{b.-end-to-end-request-processing}

For a request to https://app.team1.test:8443/api/status, the client
first asks its configured resolver for an IPv4 A record. The resolver
returns the edge address. The client opens a TCP connection to port
8443, negotiates TLS, validates the server certificate, and sends an
HTTP GET inside the encrypted session. nginx terminates TLS, chooses an
upstream and establishes a separate HTTP connection to port 3001 or
3002. The backend returns JSON and an X-Backend header; nginx forwards
the response to the client over HTTPS. DNS is not in the application
data path after resolution \cite{ref3}, \cite{ref4}.

The edge preserves application metadata with Host, X-Real-IP,
X-Forwarded-For and X-Forwarded-Proto headers. In particular,
X-Forwarded-Proto identifies the original HTTPS scheme even though the
backend hop is HTTP. The client does not need to know the backend IP
addresses. The edge can be related to a cloud application load balancer,
while the DNS node is analogous to a private managed DNS service;
neither cloud service was used in the implementation.

\subsection{Protocol and Tool
Selection}\label{c.-protocol-and-tool-selection}

The design uses IPv4 and local wireless link delivery, ICMP for
reachability tests, UDP or TCP for DNS, TCP for reliable byte-stream
transport, TLS for confidentiality and authentication, and HTTP for
request-response semantics \cite{ref3}--\cite{ref8}. Homebrew provides dnsmasq
and nginx; Python supplies the backend HTTP server; OpenSSL supports
certificate verification; curl, dig, nc, lsof and Wireshark provide
diagnostic and observational evidence.

\section{PHASE 1: BUILD AND
OBSERVE}\label{iii.-phase-1-build-and-observe}

\subsection{LAN Setup and Private
DNS}\label{a.-lan-setup-and-private-dns}

The team documented machine roles, IPv4 addressing and connectivity
using interface and ping tests. An IP address identifies a host
interface, while a destination port selects a service on that host.
Successful ICMP echo exchanges establish reachability but do not
establish that an application is listening. The active edge uses TCP
8081 for HTTP and TCP 8443 for HTTPS; these are explicit, non-default
service ports.

Mac 1 runs dnsmasq with local mappings for app.team1.test and
api.team1.test. The .test namespace is reserved for testing \cite{ref9},
avoiding the .local multicast DNS namespace. Public names are forwarded
to 8.8.8.8 and 1.1.1.1; the Phase 1 config sets cache-size=1000. Client
resolver settings were configured to query the team DNS. Direct dig
queries and public-name checks verified private resolution and
forwarding. This isolates the meaning of a successful DNS query from the
later TCP connection.

\subsection{Backend Services and Reverse
Proxy}\label{b.-backend-services-and-reverse-proxy}

Backend A and Backend B use Python's BaseHTTPRequestHandler and
HTTPServer, binding to 0.0.0.0 on ports 3001 and 3002, respectively.
Binding to all IPv4 interfaces permits LAN clients and the edge to
connect; a loopback-only binding would prevent remote access. Both
services implement GET / and GET /api/status. The home endpoint returns
HTML identifying the serving node; the status endpoint returns JSON such
as \{"backend": "A", "status": "ok"\}. An X-Backend header identifies
the actual upstream, and unknown routes return HTTP 404. No database or
shared application state is involved \cite{ref2}.

nginx uses a common upstream group named backend\_servers for both its
HTTP and HTTPS server blocks. Without another balancing directive, it
distributes requests using round robin. The Phase 1 configuration
defined max\_fails=3, fail\_timeout=10s, a 5 s connection timeout and a
60 s read timeout. Repeated requests to the same hostname yielded A and
B identifiers, showing that the edge rather than DNS distributes
individual requests. Round-robin selection is not a guarantee of strict
alternation when retries, failures or multiple workers are involved.

\subsection{TLS Termination and Certificate
Trust}\label{c.-tls-termination-and-certificate-trust}

The local CA, Team1 Local CA, issued the certificate for app.team1.test.
The server certificate and private key were installed in the nginx
certificate directory; the local CA was trusted by test clients.
Repository TLS notes record certificate validity from 29 September 2026
to 29 September 2027 and a successful certificate verification. The TLS
configuration permits TLS 1.2 and TLS 1.3. The captured connection
negotiated TLS 1.3, cipher TLS\_CHACHA20\_POLY1305\_SHA256 and ALPN
http/1.1 \cite{ref2}, \cite{ref6}.

The TCP three-way handshake (SYN, SYN-ACK, ACK) establishes a reliable
transport connection before TLS. In TLS 1.3, the ClientHello and
ServerHello establish negotiation and key exchange parameters;
authentication and Finished messages complete the secure session. Much
of the TLS 1.3 handshake after ServerHello is encrypted, so a passive
packet capture alone cannot necessarily reveal the certificate contents.
TLS protects the client-to-edge hop; traffic between nginx and the
backends remains plaintext HTTP. The private key is not included in the
GitHub submission.

\subsection{HTTP Cache Validation}\label{d.-http-cache-validation}

The root endpoint emits Cache-Control: public, max-age=30 and ETag:
"team1-home-v1". When a client sends an If-None-Match header with this
tag, the handler returns HTTP 304 Not Modified without transmitting a
new response body. A fresh cache hit can avoid an HTTP request entirely;
conditional validation still contacts the server but avoids
retransmitting an unchanged representation. The /api/status endpoint
uses Cache-Control: no-store, so health observations are not intended to
be cached \cite{ref7}, \cite{ref8}.

The following commands use GET instead of HEAD because the Python
handler implements do\_GET but not do\_HEAD. The expected first response
is HTTP 200 with cache headers, followed by HTTP 304 for the conditional
request.

curl -sS -D - -o /dev/null \textbackslash{}\\
https://app.team1.test:8443/\\
curl -sS -D - -o /dev/null \textbackslash{}\\
-H \textquotesingle If-None-Match: "team1-home-v1"\textquotesingle{}
\textbackslash{}\\
https://app.team1.test:8443/

The shared ETag is sufficient for the classroom demonstration but is not
correct representation-specific validation for differing
backend-identifying HTML bodies. A production design would use identical
representations or representation-specific validators. No nginx response
cache was configured; the experiment concerns HTTP caching semantics at
the application/client boundary.

\subsection{Packet-Level
Observation}\label{e.-packet-level-observation}

The Phase 1 packet capture, captures/mac4\_capture.pcapng, records
traffic from Mac 4 using the earlier lab addresses. Table II reproduces
the packet identifiers documented in the repository. The DNS
query/response locates the edge. The TCP handshake establishes the
client-edge connection, while ClientHello includes the Server Name
Indication (SNI) app.team1.test. TLS then protects the frontend HTTP
payload; the subsequent separate backend HTTP exchange is visible in
plaintext \cite{ref2}.

\begin{table}[!t]
\caption{DOCUMENTED PHASE 1 PACKET TRACE}
\centering\footnotesize
\begin{tabular}{@{}
  >{\raggedright\arraybackslash}p{(\columnwidth - 4\tabcolsep) * \real{0.2631}}
  >{\raggedright\arraybackslash}p{(\columnwidth - 4\tabcolsep) * \real{0.1959}}
  >{\raggedright\arraybackslash}p{(\columnwidth - 4\tabcolsep) * \real{0.5409}}@{}}
\toprule
\begin{minipage}[b]{\linewidth}\raggedright
\textbf{Packet(s)}
\end{minipage} & \begin{minipage}[b]{\linewidth}\raggedright
\textbf{Protocol}
\end{minipage} & \begin{minipage}[b]{\linewidth}\raggedright
\textbf{Observation}
\end{minipage} \\
\midrule



1439, 1463 & DNS & A query and answer: 10.7.20.19 \\
2307, 2321, 2324 & TCP & SYN, SYN-ACK, ACK to 8443 \\
2327, 2330 & TLS & ClientHello SNI; ServerHello TLS 1.3 \\
2365 & HTTP & Forwarded GET to B; X-Forwarded-Proto: https \\
2368 & HTTP & HTTP/1.0 200 OK; X-Backend: B \\
\bottomrule
\end{tabular}
\end{table}

TCP uses sequence and acknowledgment numbers, retransmission and
receive-window flow control to deliver an ordered byte stream; the
client source port is ephemeral and the destination port is fixed by the
requested service. The documented capture shows separate TCP connections
for the encrypted frontend and plaintext upstream. This is consistent
with TLS termination at the reverse proxy. Representative Wireshark
display filters are dns, tcp.port == 8443, tls, and tcp.port == 3002
\&\& http.

\subsection{Phase 1 Controlled Failure
Tests}\label{f.-phase-1-controlled-failure-tests}

The five baseline fault experiments, preserved as screenshots and
summarized in the repository, established that service health cannot be
inferred from DNS or ping alone. Table III reports the recorded
observations, not universal error codes. For example, an unreachable
upstream may produce HTTP 502 or 504 depending on whether the connection
is refused or times out.

\begin{table}[!t]
\caption{BASELINE FAILURE OBSERVATIONS}
\centering\footnotesize
\begin{tabular}{@{}
  >{\raggedright\arraybackslash}p{(\columnwidth - 2\tabcolsep) * \real{0.4796}}
  >{\raggedright\arraybackslash}p{(\columnwidth - 2\tabcolsep) * \real{0.5204}}@{}}
\toprule
\begin{minipage}[b]{\linewidth}\raggedright
\textbf{Injected condition}
\end{minipage} & \begin{minipage}[b]{\linewidth}\raggedright
\textbf{Recorded outcome}
\end{minipage} \\
\midrule



Incorrect HTTPS port (8444) & TCP connection refused \\
Backend A unavailable & Six HTTP 200 replies from B \\
Both backends unavailable & nginx HTTP 504 Gateway Time-out \\
Query sent to wrong DNS server & dig timed out \\
Wrong hostname/IP mapping & HTTPS to incorrect host timed out \\
\bottomrule
\end{tabular}
\end{table}

\section{PHASE 2: HARDEN AND
RECOVER}\label{iv.-phase-2-harden-and-recover}

\subsection{Backup DNS Resolver (Extension
A)}\label{a.-backup-dns-resolver-extension-a}

A second dnsmasq instance was installed on Mac 4. Both resolvers listen
on their respective LAN addresses and loopback, serve matching A
records, use local=/team1.test/, and forward public queries upstream.
Mac 2 lists 10.185.128.177 followed by 10.185.128.32. Mac 4 also uses
both team servers. Identical private records reduce inconsistency if the
primary is unavailable, but the two configurations are maintained
independently rather than replicated automatically.

The primary service was stopped with Homebrew while the backup, edge and
backends remained operational. In the saved 24 November log, a direct
query to the primary timed out, while the backup returned NOERROR, the
30 s A record and a measured query time of 92 ms. The system host lookup
succeeded and the HTTPS status request returned HTTP 200 from B. After
restart, three direct primary queries succeeded and an HTTPS request
returned HTTP 200 \cite{ref2}. The result establishes backup availability
and service continuity. The system lookup may have used a cached result,
so the log alone does not prove which resolver served that lookup. A
cold-cache client lookup and simultaneous DNS capture would strengthen
the proof.

\subsection{DNS TTL and Controlled Record Change (Extension
B)}\label{b.-dns-ttl-and-controlled-record-change-extension-b}

Both resolvers were configured with local-ttl=30, causing locally
defined A records to advertise a 30 s lifetime. The saved validation log
confirms this TTL for app.team1.test on both resolvers and
api.team1.test on the primary. During an earlier interactive experiment,
the primary api.team1.test mapping was changed from 10.185.128.147 to
10.185.128.163. Repeated macOS system lookups on Mac 4 showed the old
address until 02:53:46 and the new address from 02:53:48. The original
mapping and dual-DNS settings were subsequently restored.

The observed transition is consistent with caching, but the
record-change instant was not preserved alongside the lookup sequence.
Consequently, the exact expiration interval cannot be inferred. A
rigorous repeat would timestamp the authoritative update, query the
client at regular intervals, record the first new answer, and separately
demonstrate a manual cache flush. The 30 s advertised TTL is not a
guarantee that all clients or existing TCP sessions migrate within 30 s.

\subsection{Service Isolation Using PF (Extension
C)}\label{c.-service-isolation-using-pf-extension-c}

macOS Packet Filter (PF) rules were loaded on Mac 3 and Mac 4 using a
dedicated team1\_backend anchor. On each backend, a quick pass rule
permitted TCP traffic from Mac 2 to its private service port, while a
subsequent quick block-drop rule denied other remote sources on en0. The
rules were stateful and limited to the relevant service port, rather
than blocking all traffic. The original /etc/pf.conf files were backed
up before changes.

\# Mac 3: saved rule intent\\
pass in quick on en0 inet proto tcp \textbackslash{}\\
from 10.185.128.147 to 10.185.128.163 \textbackslash{}\\
port 3001 keep state\\
block drop in quick on en0 inet proto tcp \textbackslash{}\\
from any to 10.185.128.163 port 3001

The corresponding Backend B rule uses destination 10.185.128.32 and port
3002. Tests using nc from Mac 2 succeeded for both ports; the same
direct tests from Mac 1 timed out. Edge load balancing remained
operational. During standby migration, a temporary exception on Mac 4
permitted the standby edge at 10.185.128.163 to reach B; local A access
used 127.0.0.1. The exception was removed after restoring the original
edge.

At the end of the demonstration, the original PF configuration files
were restored, project anchor rules were cleared and the activation
references released. PF reported Disabled on both backend Macs and Mac 1
could again connect directly to both ports. This demonstrates
access-control behavior and safe rollback, not persistent firewall
enforcement in the final operating state. The rules implement IPv4
source-based filtering, not cryptographic identity or complete host
hardening.

\subsection{Backend High Availability (Extension
D)}\label{d.-backend-high-availability-extension-d}

The active upstream configuration was strengthened by reducing
max\_fails from 3 to 1 within fail\_timeout=10s. The connection timeout
was reduced from 5 s to 2 s and the read timeout from 60 s to 30 s.
proxy\_next\_upstream specifies error, timeout, HTTP 502, 503 and 504;
proxy\_next\_upstream\_tries=2 limits attempts. This is passive failure
handling: nginx observes unsuccessful requests, temporarily avoids an
unavailable upstream and retries an eligible request on the other
upstream. It is not an independent periodic active health check.

upstream backend\_servers \{\\
server 10.185.128.163:3001\\
max\_fails=1 fail\_timeout=10s;\\
server 10.185.128.32:3002\\
max\_fails=1 fail\_timeout=10s;\\
\}

Backend A was stopped and ten successive HTTPS status requests returned
HTTP 200 from Backend B. After A was restarted, both identifiers
appeared again. A later saved healthy-state validation on 24 November
contained ten HTTP 200 responses, four served by A and six by B (Fig.
2). The outage result is a retrospective account; the healthy-state
distribution is recorded in a contemporaneous log. The ten-request
sample demonstrates functionality but cannot establish an availability
percentage, a guaranteed recovery time or throughput under load.

\begin{figure}[!t]\centering
\includegraphics[width=\columnwidth]{ieee_tex_assets/media/image2.png}
\caption{Observed backend distribution from the saved 03:53:33 IST
validation log on 24 November 2026.}
\end{figure}

\subsection{Controlled Edge Migration (Extension
E)}\label{e.-controlled-edge-migration-extension-e}

A standby nginx was configured on Mac 3 at 10.185.128.163:8443. Its
upstream group reaches Backend A through loopback and Backend B over the
LAN. The project TLS certificate and key were provisioned separately,
and certificate validation against the project CA succeeded. Before
changing DNS, six HTTPS requests using curl -\/-resolve to the standby
address returned HTTP 200 and reached both backends. The original edge
continued running during this controlled migration.

Both DNS servers were updated so that app.team1.test resolved to the
standby IP; api.team1.test remained pointed at the active edge. After
restarting the DNS services, both returned the new address with a 30 s
TTL. The client system resolver returned 10.185.128.163, and a normal
HTTPS request without -\/-resolve returned HTTP 200 from nginx 1.31.9
and Backend A. The original mapping was then restored to 10.185.128.147.
A normal HTTPS request returned HTTP 200 from nginx 1.31.6 and Backend
B. A later saved validation log confirms six successful requests, three
from each backend, on the restored edge.

This test shows that a hostname can be reassigned to a functioning
standby edge without changing the client URL. The migration is manual,
not automatic failure detection or instantaneous failover. The test did
not record a simultaneous old-cache/new-lookup comparison at the cutover
instant. Existing connections and cached DNS records can continue to
reach the old edge until expiry or reconnection. After restoration,
standby nginx and temporary SSH access were stopped; the configuration
and certificate files remained on Mac 3.

\subsection{Fault Diagnosis and Recovery (Extension
F)}\label{f.-fault-diagnosis-and-recovery-extension-f}

The team rehearsed the faculty-injected troubleshooting procedure by
stopping active nginx. From Mac 1, a direct DNS query returned the
correct IP, ICMP ping succeeded with three of three replies, and a TCP
connection to 8443 was refused. A normal curl request additionally
failed because the system DNS resolver on Mac 1 was configured as
8.8.8.8. A curl -\/-resolve test bypassed DNS while preserving the
hostname, confirming that the TCP failure was independent of the
name-resolution error.

Mac 1 was reconfigured to use its loopback resolver and the backup at
10.185.128.32, restoring system resolution. On Mac 2, nginx -t passed,
nginx was started and lsof verified that port 8443 was listening.
Subsequent HTTPS requests returned HTTP 200 and six requests reached
both backends. The fault was therefore isolated at two layers: incorrect
client resolver configuration and an absent TCP listener. This was a
team rehearsal, not evidence that the actual faculty-injected assessment
had already taken place.

\section{RESULTS AND EVALUATION}\label{v.-results-and-evaluation}

The experimental outcomes are summarized in Table IV. Direct log
evidence and retrospective observations are differentiated because they
have different evidentiary strength. All observations relate to short,
controlled experiments on the team's LAN, with no synthetic load test or
repeated timing campaign.

\begin{table}[!t]
\caption{PHASE 2 EXPERIMENT SUMMARY}
\centering\footnotesize
\begin{tabular}{@{}
  >{\raggedright\arraybackslash}p{(\columnwidth - 2\tabcolsep) * \real{0.2865}}
  >{\raggedright\arraybackslash}p{(\columnwidth - 2\tabcolsep) * \real{0.7135}}@{}}
\toprule
\begin{minipage}[b]{\linewidth}\raggedright
\textbf{Extension}
\end{minipage} & \begin{minipage}[b]{\linewidth}\raggedright
\textbf{Observed result / evidence}
\end{minipage} \\
\midrule



A: DNS redundancy & Primary timeout; backup NOERROR; HTTPS 200. Saved
logs. \\
B: TTL & 30 s A-record TTL in saved log; old-to-new client lookup
observed earlier. \\
C: Isolation & Edge access allowed; Mac 1 timed out; rollback verified.
Written account. \\
D: Backend HA & 10/10 HTTP 200 from B during A outage (written account);
healthy log 4 A, 6 B. \\
E: Edge migration & New IP .163 and HTTPS 200; restoration .147 and 6/6
successes. Summary + recovery log. \\
F: Diagnosis & DNS and TCP faults isolated; nginx restored; six
successful requests. Summary + recovery log. \\
\bottomrule
\end{tabular}
\end{table}

The redundancy tests show that the architecture can tolerate selected
service failures. DNS resilience addresses the directory function;
upstream failover addresses the application execution path. These are
independent. If the active edge itself fails, DNS may still resolve
successfully but clients cannot establish HTTPS until an alternative
edge is activated. The active edge remains a single point of failure in
the default restored state.

No latency benchmark is claimed from the isolated 92 ms DNS response.
The observed request counts are validation samples rather than
statistical estimates of long-term reliability. Likewise, a successful
TLS request demonstrates the observed hostname and trust chain, not
encryption of the internal backend segment.

\section{REPRODUCIBILITY AND
OPERATIONS}\label{vi.-reproducibility-and-operations}

\subsection{Baseline Service
Startup}\label{a.-baseline-service-startup}

The repository snapshot at commit 2ae489c contains the Phase 1 files,
six Phase 2 configuration snapshots, and evidence folders \cite{ref2}.
Before a repeat demonstration, verify the current Wi-Fi IPv4 addresses
and update DNS records, nginx upstreams and PF rules together. All
devices must remain on a network that permits peer communication. The
repository snapshots use demonstration-specific paths and addresses and
must be adapted before deployment.

On Macs 1 and 4, validate the installed dnsmasq configuration and start
the service if it is not already running:

sudo /opt/homebrew/opt/dnsmasq/sbin/dnsmasq \textbackslash{}\\
-\/-test -\/-conf-file=/opt/homebrew/etc/dnsmasq.conf\\
sudo brew services start dnsmasq

On Mac 3, launch Backend A in a dedicated Terminal; on Mac 4, substitute
backend-b. Stop the Python foreground process with Ctrl-C only during
controlled tests.

cd \textasciitilde/Team\_1.1.1.1/backends/backend-a\\
python3 server.py

On Mac 2, check nginx syntax and launch it if stopped. For an already
running instance, use nginx -s reload after a successful syntax test.
Ensure the TLS files are provisioned locally and the project CA is
trusted.

nginx -t \&\& nginx\\
lsof -nP -iTCP:8443 -sTCP:LISTEN

\subsection{Layer-by-Layer
Verification}\label{b.-layer-by-layer-verification}

The diagnostic sequence starts with network identity and DNS, then TCP,
TLS and HTTP. The following commands reproduce the primary DNS query,
system lookup, listener reachability and trusted application request
from an appropriately configured client:

ipconfig getifaddr en0\\
networksetup -getdnsservers "Wi-Fi"\\
dig @10.185.128.177 app.team1.test\\
dscacheutil -q host -a name app.team1.test\\
nc -G 3 -vz 10.185.128.147 8443\\
curl -v -\/-noproxy \textquotesingle*\textquotesingle{}
\textbackslash{}\\
https://app.team1.test:8443/api/status

A DNS timeout indicates that the queried resolver did not respond; an
incorrect A record indicates incorrect destination selection; TCP
refusal suggests no service listening or an explicit reject; a TLS
verification failure points to certificate, hostname or trust; HTTP
502/504 indicates a reachable edge with an upstream problem. A direct
curl -\/-resolve test can separate resolver faults from transport or TLS
faults without bypassing certificate validation.

\subsection{Evidence and Configuration
Management}\label{c.-evidence-and-configuration-management}

The team pushed its Phase 2 work to the main branch as commit 2ae489c
(21 files changed, 522 insertions), after checking staged files and
removing carriage-return formatting from captured HTTP headers.
README.md, MANIFEST.md and PHASE2.md index the work. The
configuration/phase2 directory holds primary and backup DNS
configurations, active and standby nginx configurations, and both
backend PF rules. Phase 1 screenshots and the packet capture are
retained; Phase 2 evidence is grouped by extension under
evidence/09-phase2. No TLS private key is included in the committed
files.

\section{LIMITATIONS AND FUTURE
WORK}\label{vii.-limitations-and-future-work}

First, the current DNS redundancy uses manually maintained records;
automated configuration synchronization and health-aware DNS would
reduce divergence. Second, a second active edge or a managed virtual IP
with automated health checks would reduce the edge single point of
failure. Third, encryption or mutual TLS on the upstream links would
extend confidentiality and authentication beyond the frontend. Fourth,
permanent least-privilege firewall policy would require a
deployment-specific security design and approval, instead of temporary
laboratory PF rules.

Fifth, HTTP validators should reflect actual representation bytes and
the Python server could be extended to support HEAD and concurrent
requests. Finally, the evidence could be improved by capturing the
automatic DNS resolver selection after cache expiry, recording
timestamped TTL changes and explicit cache flushes, capturing both
cutover paths concurrently, and running repeatable timed fault and load
tests. These improvements are proposed, not claimed as implemented.

\section{CONCLUSION}\label{viii.-conclusion}

Team 1.1.1.1 implemented a complete local request path from private DNS
resolution through TCP, TLS and nginx to two Python HTTP backends. Phase
1 demonstrated baseline connectivity, secure reverse proxying, load
balancing, conditional HTTP responses and packet-level analysis. Phase 2
extended the same platform with backup DNS, short TTLs, controlled
backend isolation, passive upstream recovery, standby edge migration and
systematic diagnosis. The tests demonstrate continuity under selected
DNS and backend failures and successful restoration following deliberate
faults. The work also establishes the limits of this design: an active
edge must be reachable, DNS migration is manual, and firewall
enforcement was intentionally rolled back. The resulting platform is a
reproducible educational model of network-service architecture and
resilience.

\begin{thebibliography}{9}
\bibitem{ref1} Computer Networks Course Project, ``Private Network Service
Platform: A Two-Phase, Team-Based Local Networking Project,'' course
handout, 2026.

\bibitem{ref2} Team 1.1.1.1, ``Computer Networks Project,'' GitHub repository,
commit 2ae489c, Nov. 24, 2026. {[}Online{]}. Available:
https://github.com/shriharikn17/Team\_1.1.1.1

\bibitem{ref3} P. Mockapetris, ``Domain names---concepts and facilities,'' RFC
1034, Nov. 1987. doi: 10.17487/RFC1034.

\bibitem{ref4} P. Mockapetris, ``Domain names---implementation and
specification,'' RFC 1035, Nov. 1987. doi: 10.17487/RFC1035.

\bibitem{ref5} W. Eddy, Ed., ``Transmission Control Protocol (TCP),'' RFC 9293,
Aug. 2022. doi: 10.17487/RFC9293.

\bibitem{ref6} E. Rescorla, ``The Transport Layer Security (TLS) Protocol
Version 1.3,'' RFC 8446, Aug. 2018. doi: 10.17487/RFC8446.

\bibitem{ref7} R. Fielding, M. Nottingham, and J. Reschke, ``HTTP Semantics,''
RFC 9110, Jun. 2022. doi: 10.17487/RFC9110.

\bibitem{ref8} R. Fielding, M. Nottingham, and J. Reschke, ``HTTP Caching,''
RFC 9111, Jun. 2022. doi: 10.17487/RFC9111.

\bibitem{ref9} D. Eastlake 3rd and A. Panitz, ``Reserved Top Level DNS Names,''
RFC 2606, Jun. 1999. doi: 10.17487/RFC2606.

\end{thebibliography}
\end{document}
